\documentclass[aspectratio=169]{beamer}
\mode<presentation> {
\usetheme{Madrid}
\setbeamertemplate{navigation symbols}{} % Remove os símbolos de navegação para um visual mais limpo
}

\usepackage{graphicx}
\usepackage{booktabs}
\usepackage[utf8]{inputenc}
\usepackage{color}

%----------------------------------------------------------------------------------------
%	TITLE PAGE
%----------------------------------------------------------------------------------------
\title[UnBlabla Car]{UnBlabla Car}
\subtitle{Mobilidade Inteligente, Atrito Zero e Networking na UnB}

\author[Equipe]{Antonio Paulo \and Lucas Drumond \and Heitor Estevam \and Luis Bessa \and Gabriel Santana}

\institute[UnB]{
    Sistemas de Informação - CIC0101 \\ 
    Universidade de Brasília
}
\date{\today}

\begin{document}

\begin{frame}
\titlepage
\end{frame}

%------------------------------------------------
% APRESENTADOR 1: O ESTRATEGISTA
%------------------------------------------------
\section{O Problema}
\begin{frame}
\frametitle{Problema de Mobilidade}
Existe uma ineficiência massiva no transporte diário até o campus Darcy Ribeiro:

\vspace{0.5cm}
\begin{itemize}
    \item \textbf{Desperdício de Capital:} Veículos circulam diariamente com assentos vazios, caracterizando subutilização de recursos.
    \item \textbf{Sobrecarga de Infraestrutura:} Estudantes dependem de transporte público sujeito a atrasos e superlotação.
    \item \textbf{Desconexão Orgânica:} O isolamento no trajeto impede o networking acadêmico entre diferentes semestres e cursos.
\end{itemize}

\vspace{0.5cm}
\textbf{Nossa Missão:} Conectar a demanda com a oferta de forma segura e eficiente.
\end{frame}

\begin{frame}
\frametitle{A Solução: UnBlabla Car}
\begin{columns}[c]
\column{.5\textwidth} % Coluna de texto ajustada para 50%
\textbf{O que o sistema É:}
\begin{itemize}
    \item Uma plataforma peer-to-peer focada exclusivamente na comunidade da UnB.
    \item Um mecanismo automatizado de divisão de custos para incentivar a congregação acadêmica.
\end{itemize}
\vspace{0.3cm}
\textbf{O que o sistema NÃO É:}
\begin{itemize}
    \item Não é um serviço de transporte comercial.
    \item Não visa a obtenção de lucro pelos motoristas ou pela própria plataforma.
\end{itemize}

\column{.45\textwidth} % Coluna de imagem expandida
\centering
% IMAGEM: HOME / MAPA
\includegraphics[height=0.85\textheight]{home_mapa.jpg}
\end{columns}
\end{frame}

%------------------------------------------------
% APRESENTADOR 2: O ANALISTA DE SEGURANÇA
%------------------------------------------------
\section{Autenticação}
\begin{frame}
\frametitle{Segurança e Acesso Restrito}
\begin{columns}[c]
\column{.45\textwidth}
\centering
% IMAGEM: LOGIN
\includegraphics[height=0.85\textheight]{login.jpg}

\column{.5\textwidth}
\textbf{Controle de Acesso (Épico 1)}
\vspace{0.3cm}
\begin{itemize}
    \item \textbf{Filtro Institucional:} O sistema exige vínculo comprovado. O acesso é estritamente restrito ao provedor oficial da UnB: \texttt{@aluno.unb.br}.
    \item \textbf{Mitigação de Riscos:} Elimina perfis falsos e usuários externos à universidade, garantindo a integridade e segurança dos passageiros.
    \item \textbf{Privacidade por Design:} Não há senhas expostas ou dados pessoais públicos na camada de entrada.
\end{itemize}
\end{columns}
\end{frame}

%------------------------------------------------
% APRESENTADOR 3: O ARQUITETO DE LOGÍSTICA
%------------------------------------------------
\section{Matching}
\begin{frame}
\frametitle{Topologia de Oferta (Motorista)}
\begin{columns}[c]
\column{.5\textwidth}
\textbf{Publicação de Rotas (Épico 2)}
\vspace{0.3cm}
\begin{itemize}
    \item \textbf{Fluxo Intuitivo:} O motorista compartilha um trajeto que \textit{já faria}, informando pontos de saída, destino, data e horário.
    \item \textbf{Controle de Capacidade:} Definição exata de vagas disponíveis no veículo.
    \item \textbf{Regras de Engajamento:} Espaço para comunicação de observações operacionais (ex: pontualidade, ponto exato no bloco).
\end{itemize}

\column{.45\textwidth}
\centering
% IMAGEM: OFERECER CARONA
\includegraphics[height=0.85\textheight]{oferecer.jpg}
\end{columns}
\end{frame}

\begin{frame}
\frametitle{Confirmação e Disponibilidade}
\begin{columns}[c]
\column{.45\textwidth}
\centering
% IMAGEM NOVA: FEEDBACK DE SUCESSO
\includegraphics[height=0.85\textheight]{07-carona-publicada.png}

\column{.5\textwidth}
\textbf{Máquina de Estados: Aberta}
\vspace{0.3cm}
\begin{itemize}
    \item \textbf{Feedback Imediato:} Confirmação visual de que a carona foi adicionada ao sistema.
    \item \textbf{Visibilidade:} A rota passa a constar no algoritmo de busca para pareamento com passageiros compatíveis.
    \item \textbf{Notificações:} O motorista entra em estado de escuta, aguardando os alertas do sistema para solicitações de vagas.
\end{itemize}
\end{columns}
\end{frame}

%------------------------------------------------
% APRESENTADOR 4: O NEGOCIADOR
%------------------------------------------------
\section{Logística}
\begin{frame}
\frametitle{Atrito Zero: Detalhamento e Demanda}
\begin{columns}[c]
\column{.5\textwidth}
\textbf{Transparência Financeira (Épico 3)}
\vspace{0.3cm}
\begin{itemize}
    \item \textbf{Cálculo Auditável:} O sistema decompõe o custo (trecho em comum + desvio da rota), removendo qualquer tensão social sobre o rateio.
    \item \textbf{Sem Pagamento no App:} A plataforma atua puramente como canal entre os passageiros e motoristas, organizando a carona.
    \item \textbf{Solicitação:} O passageiro aciona o pedido e aguarda o aceite para iniciar o protocolo de comunicação (WhatsApp).
\end{itemize}

\column{.45\textwidth}
\centering
% IMAGEM NOVA: DETALHES DA CARONA COM CÁLCULO
\includegraphics[height=0.85\textheight]{04-detalhes-carona-solicitacao.png}
\end{columns}
\end{frame}

\begin{frame}
\frametitle{Gestão de Vagas: Visão Dual}
\begin{columns}[c]
\column{.45\textwidth}
\centering
% IMAGEM NOVA: VISÃO PASSAGEIRO
\includegraphics[height=0.85\textheight]{05-minhas-viagens-passageiro.png}

\column{.5\textwidth}
\textbf{Comando e Controle Simultâneo}
\vspace{0.3cm}
\begin{itemize}
    \item \textbf{Passageiro:} Monitoramento do status "Pendente" com instruções claras sobre o cronograma com horários e locais de embarque/desembarque.
    \item \textbf{Motorista:} Soberania total. Visualização do perfil do solicitante, curso e avaliação, com decisão de aceitação ou recusa.
    \item \textbf{Comunicação:} Os canais de comunicação são liberados para conversa entre os usuários para alinhamento das expectativas.
\end{itemize}
\end{columns}
\end{frame}

% Adicionando a imagem do motorista lado a lado ou em sequência rápida
\begin{frame}
\frametitle{Gestão de Vagas: Decisão do Motorista}
\begin{columns}[c]
\column{.5\textwidth}
\textbf{Autorização de Embarque}
\vspace{0.3cm}
\begin{itemize}
    \item \textbf{Filtro de Confiança:} O motorista observa os detalhes do passageiro (Nota e Curso) antes de permitir a entrada no veículo.
    \item \textbf{Atualização de Status:} Ao aceitar, o sistema consome uma vaga disponível da oferta original e notifica o passageiro.
\end{itemize}

\column{.45\textwidth}
\centering
% IMAGEM NOVA: VISÃO MOTORISTA
\includegraphics[height=0.85\textheight]{08-minhas-caronas-pedidos.png}
\end{columns}
\end{frame}

%------------------------------------------------
% APRESENTADOR 5: O GESTOR DE COMUNIDADE
%------------------------------------------------
\section{Reputação}
\begin{frame}
\frametitle{Confiança, Gamificação e Sustentabilidade}
\begin{columns}[c]
\column{.45\textwidth}
\centering
% IMAGEM: PERFIL
\includegraphics[height=0.85\textheight]{perfil.jpg}

\column{.5\textwidth}
\textbf{O Futuro do Ecossistema (Épico 4)}
\vspace{0.3cm}
\begin{itemize}
    \item \textbf{Métricas de Confiança:} Visualização do perfil, com sistema de notas atribuídas (1 a 5 estrelas) e histórico de viagens do usuário.
    \item \textbf{Incentivos ao Usuário:} Gamificação através do cálculo de Kg de CO2 economizados ao preencher assentos que estariam vazios e recebimento de medalhas/troféus ao realizar viagens e compartilhar caronas.
\end{itemize}
\end{columns}
\end{frame}

%------------------------------------------------
\begin{frame}
\Large{\centerline{Obrigado!}}
\Large{\centerline{Dúvidas?}}

\vspace{1cm}
\centering{
    \color{blue}Sistemas de Informação - UnB
}
\end{frame}

\end{document}